\documentclass[10pt]{article}

\usepackage[margin=0.52in]{geometry}
\usepackage[T1]{fontenc}
\usepackage{array,booktabs,hyperref}
\usepackage{fancyhdr,float}

\newcommand{\course}{CS 571}
\newcommand{\assignment}{Homework 2 -- Entity and Relation Extraction}
\newcommand{\studentname}{Shenghan Wen}
\newcommand{\studentid}{2823256}

\hypersetup{colorlinks=true,urlcolor=blue}
\pagestyle{fancy}
\fancyhf{}
\lhead{\course\ -- Homework 2}
\rhead{\studentname}
\cfoot{\thepage}
\setlength{\headheight}{14pt}
\setlength{\parindent}{0pt}
\setlength{\parskip}{0.24em}
\setlength{\tabcolsep}{3.5pt}
\renewcommand{\arraystretch}{1.08}
\title{\vspace{-1.85cm}\large\course\\[0.12em]\normalsize\assignment}
\author{\small\studentname\quad (\studentid)}
\date{}

\begin{document}
\fontsize{9}{10.4}\selectfont
\maketitle
\thispagestyle{fancy}
\vspace{-1.55em}

This homework is finished with the assistance of Codex (GPT-5.6 Sol).

\section*{I. Supervised Entity Tagging}

\textbf{Task and protocol.}
I formulate NER as token-level BIO classification over PER, ORG, LOC, TIME,
NUM, and MISC.  A prediction is correct only when sentence, start, end, and type
exactly match a gold mention.  Model and ensemble choices use the 500-document
dev split; test is evaluated once after locking both stages.

\textbf{Model and features.}
The primary tagger is \texttt{bert-base-cased}, first trained on the 10,000
weakly labeled documents and then fine-tuned on human labels.  Contextual
WordPiece features encode identity, subword morphology, capitalization,
punctuation, and bidirectional sentence context.  The first WordPiece receives
the source token's BIO label; special and later pieces are masked.  A supervised
\texttt{distilbert-base-cased} companion provides an independent vote.  Dev
selection keeps every primary PER span and exact agreement between both models
for ORG, LOC, TIME, NUM, and MISC.  Orphan I-tags become B-tags.  Mentions are
grouped iff type and case-folded, whitespace-normalized surface string match.

\textbf{Selected hyperparameters.}
The primary uses one weak epoch and two human epochs, batch size 32, AdamW at
$3\times10^{-5}$, weight decay .01, 10\% warmup, linear decay, clipping 1.0,
maximum length 512, O-label weight .60, and seed 571.  Length-grouped batches
reduce padding.  The companion uses two human epochs at $5\times10^{-5}$.

\textbf{Final span-level results.} The locked system gains .0197 test micro-F1
over the supervised baseline (.8747 to .8944), approximately two points.

\begin{table}[H]
\centering\scriptsize
\begin{tabular}{@{}l c c c c c c@{}}
\toprule
& \multicolumn{3}{c}{Dev} & \multicolumn{3}{c}{Test} \\
\cmidrule(lr){2-4}\cmidrule(lr){5-7}
Type & P & R & F1 & P & R & F1 \\
\midrule
PER  & .9493 & .9630 & .9561 & .9330 & .9512 & .9420 \\
ORG  & .8770 & .8053 & .8397 & .8814 & .7943 & .8356 \\
LOC  & .9442 & .9153 & .9295 & .9484 & .9093 & .9285 \\
TIME & .9539 & .9352 & .9445 & .9376 & .9212 & .9293 \\
NUM  & .8081 & .8318 & .8198 & .8229 & .8571 & .8397 \\
MISC & .8363 & .7215 & .7747 & .8427 & .7553 & .7966 \\
\midrule
\textbf{Micro} & \textbf{.9151} & \textbf{.8772} & \textbf{.8958}
               & \textbf{.9129} & \textbf{.8766} & \textbf{.8944} \\
\bottomrule
\end{tabular}
\par\vspace{0.15em}{\small\textbf{Table 1.} Exact mention span-and-type results for the final distant-supervision ensemble.}
\end{table}
\vspace{-0.6em}

\textbf{Error analysis.}
On test, 796 paired near misses contain 326 type-only, 368 boundary-only, and
102 combined errors; another 802 gold mentions are missed and 287 predictions
are spurious.  Boundary errors occur in 59.0\% of paired misses (285 too narrow,
184 too wide).  MISC$\rightarrow$ORG is the largest type confusion (61), and
MISC also has the most complete misses (260).  The precision-oriented ensemble
removes many spurious spans but sacrifices recall, especially for heterogeneous
MISC names.  Span decoding with learned boundary scores is the clearest next step.

\newpage
\thispagestyle{fancy}
\fontsize{9}{10.4}\selectfont
\section*{II. Relation Extraction}

\textbf{Two-stage protocol.}
The system is a strict cascade: NER is selected and frozen on dev; RE is then
selected using full-pipeline dev F1; both are finally evaluated once on test.

\textbf{Entity construction and candidate pairs.}
The entity construction rule is the one used in Part I: repair orphaned I-tags,
convert every BIO segment to a mention, then group mentions if and only if they
have the same predicted type and identical case-folded, whitespace-normalized
surface strings.  The groups are written in \texttt{vertexSet}; their list
positions become entity indices.  For a document with $n$ entities, relation
inference considers all $n(n-1)$ ordered pairs $(h,t)$ with $h\ne t$.  Direction
is therefore explicit.  A pair is assigned one of the 96 Wikidata relations or
\texttt{no\_relation}; self-pairs and cross-document pairs are never candidates.

\textbf{Feature set.}
For each ordered pair, sparse categorical features encode (1) head type,
tail type, and their ordered combination; (2) normalized head/tail names, name
tokens, and the ordered name pair; (3) mention-count buckets; (4) whether either
name occurs in the document title; and (5) the closest pair of mentions:
sentence distance, same-sentence status, argument order, token-distance bucket,
up to 12 tokens between the mentions, and the immediate left/right token of
each argument.  Features are mapped to $2^{19}=524{,}288$ nonnegative hashed
dimensions, which controls memory while retaining lexical conjunctions.

\textbf{Classifier and class imbalance.}
The classifier is an averaged multiclass linear model trained by SGD with
logistic loss and L2 regularization.  Settings are $\alpha=2\times10^{-6}$,
25 epochs, tolerance $10^{-3}$, and seed 571.  All positives are retained;
unrelated pairs are sampled at five per positive.  Frequency skew receives
$\min(8,(n_{\max}/n_y)^{0.2})$ weighting.  The dev-selected decoder emits at
most two relations whose probability is at least .90 times \texttt{no\_relation},
recovering simultaneous relations missed by the earlier argmax decoder.

\textbf{Output and evaluation.}
For every document, the predictor preserves its sentences and predicted
\texttt{vertexSet}, removes any input relation labels, and writes a new
\texttt{labels} list containing dictionaries \texttt{\{h,t,r\}}.  Both final
dev and test files pass \texttt{check\_format.py --re} with zero errors and are
accepted by \texttt{re\_eval.py}.  In gold-entity mode the evaluator uses entity
indices directly.  In pipeline mode it aligns a predicted entity to a gold
entity through exact shared mention spans; a triple involving an unaligned
entity is a false positive.

\begin{table}[H]
\centering\small
\begin{tabular}{@{}lllrrrrrr@{}}
\toprule
Split & Training & Entities & P & R & F1 & TP & FP & FN \\
\midrule
Dev  & supervised & Pipeline & .3631 & .2918 & .3236 & 5044 & 8847 & 12240 \\
Dev  & + distant  & Pipeline & .3850 & .3856 & \textbf{.3853} & 6664 & 10646 & 10620 \\
Test & supervised & Pipeline & .3620 & .2965 & .3260 & 5173 & 9117 & 12275 \\
Test & + distant  & Pipeline & .3842 & .3957 & \textbf{.3899} & 6904 & 11066 & 10544 \\
Test & + distant  & Gold & .4839 & .4629 & .4732 & 8077 & 8614 & 9371 \\
\bottomrule
\end{tabular}
\par\vspace{0.2em}{\small\textbf{Table 2.} Relation micro scores.  ``Pipeline''
uses this system's own NER entities; ``Gold'' isolates the relation classifier.
``+ distant'' uses the final frozen NER and dev-selected weak model/decoder.}
\end{table}

The final test pipeline gains .0639 absolute F1 over supervised training.  It
predicts 17,970 triples against 17,448 gold.  Predicted entities reach 90.7\%
of gold entities and make 89.5\% of gold triples reachable.

\subsection*{Relation error analysis and pipeline gap}
Replacing predicted entities with gold increases test precision from .3842 to
.4839, recall from .3957 to .4629, and F1 from .3899 to .4732: an .0833 gap.
Of 17,970 pipeline triples, 4,186 (23.3\%) use
an entity that cannot be aligned to gold and therefore become automatic false
positives.  Gold entities recover 1,173 correct triples, yet gold-entity recall
is only .4629, so relation classification still causes most false negatives.

Errors are concentrated in frequent and semantically adjacent relations.
\texttt{P131} and \texttt{P17} improve to F1 .5096 and .4734, but P17 produces
2,664 false positives.  \texttt{P27} recall is .7606 with 1,871 false positives,
showing weak-label overgeneralization.  Part relations remain hard
(\texttt{P527} F1 .1777; \texttt{P361} .1687), and rare labels receive few or no
predictions.  Calibrated relation-specific thresholds and entity confidence
would directly target these errors.

\newpage
\thispagestyle{fancy}
\fontsize{9}{10.4}\selectfont
\section*{III. Extra Credit: Distant Supervision}

\textbf{Database and label construction.}
I select Wikidata because DocRED's 96 predicates are already Wikidata property
IDs.  Wikidata stores item--property--value statements \cite{wikidata}.
Following DocRED \cite{docred}, mentions are linked to Wikidata items, identical
KB IDs are merged, and item pairs receive their Wikidata properties.  For
reproducibility, \texttt{prepare\_distant\_supervision.py} matches title plus
complete tokenized text to the official \texttt{train\_distant.json}, rather
than issuing time-varying live queries.  All 10,000 supplied documents match,
restoring 192,955 weak entity clusters and 147,974 weak relation labels.

\textbf{Use in NER.}
The BERT primary receives one weak epoch before two human epochs with a reset
optimizer.  The supervised DistilBERT agreement filter is selected only on dev.
This changes dev/test F1 from .8771/.8747 to .8958/.8944: +.0187/+.0197,
approximately two points, while raising test precision from .8597 to .9129.

\textbf{Use in RE.}
All weak pair rows are included with source weight .25 versus 1.0 for human
rows.  On the frozen final NER dev output, at most two relations and probability
ratio .90 are selected.  Test F1 rises from .3260 to .3899 (+.0639).

\begin{table}[H]
\centering\small
\begin{tabular}{@{}llrrr@{}}
\toprule
Stage & Split & Supervised F1 & + distant F1 & Gain \\
\midrule
NER & Dev  & .8771 & .8958 & +.0187 \\
NER & Test & .8747 & .8944 & +.0197 \\
RE pipeline & Dev  & .3236 & .3853 & +.0617 \\
RE pipeline & Test & .3260 & .3899 & +.0639 \\
\bottomrule
\end{tabular}
\par\vspace{0.15em}{\small\textbf{Table 3.} Absolute micro-F1 improvement from the final distant-supervision pipeline.}
\end{table}

\textbf{Reproducibility and output QA.}
\texttt{make train-ner}, \texttt{make predict-ner SPLIT=dev}, and
\texttt{make eval-ner SPLIT=dev} provide one-command NER operations; RE has the
parallel three targets.  The four final JSON files pass \texttt{check\_format.py}
with zero errors and both evaluators.  I also opened NER and RE test outputs in
the supplied annotation viewer: all 500 documents render, with 12,434 mentions,
six entity types, and 17,970 predicted triples across 91 relations.

\begin{thebibliography}{2}\small
\bibitem{docred} Y. Yao et al.  ``DocRED: A Large-Scale Document-Level Relation
Extraction Dataset.'' ACL 2019.  \url{https://aclanthology.org/P19-1074/}
\bibitem{wikidata} Wikimedia Foundation. ``Wikidata: Introduction'' and
``Wikidata: Data model.'' Accessed October 5, 2026.
\url{https://www.wikidata.org/wiki/Wikidata:Introduction};
\url{https://www.wikidata.org/wiki/Help:Data_model}
\end{thebibliography}

\end{document}
